\section{Results}
\label{sec:results}

\subsection{Baseline difference-in-differences}
\label{sec:res-baseline}

Table~\ref{tab:stage1_composite} reports the baseline single-outcome estimates for the official MDCH flag, the two composite outcomes, and food insecurity, under the Simple, Adjusted and Extended specifications of Equation~\ref{eq:baseline-did}. The official MDCH flag falls by 6.2 percentage points under the Simple specification and by 7.6 percentage points once the six CASE controls are added -- the coefficient becomes \emph{more}, not less, significant on adjustment, even though the added controls explain noticeably less of the outcome variance than a richer, ad-hoc covariate set considered and rejected earlier in the project, which is a useful indication that the headline result is not an artefact of a kitchen-sink specification. Any deprivation falls significantly under every specification ($-0.081$ to $-0.079$ percentage points); severe deprivation shows no significant change under any specification. The Extended specification, which adds the child's own age on top of the CASE controls, barely moves any coefficient -- a useful robustness confirmation that the result is not sensitive to whether the child's own age is controlled for.

\input{Tables/stage1_composite.tex}

Table~\ref{tab:stage1_case_exact} reports the exact replication of \citet{andersen2025}'s own Table 3 specification, which adds the standalone \texttt{Post} main-effect term omitted from Equation~\ref{eq:baseline-did} (Section~\ref{sec:baseline-did}) and reports the full covariate vector rather than only the treatment interaction. The Scotland$\times$Post coefficient of $-0.079$ for child material deprivation and $-0.068$ for food insecurity are both close to the corresponding Adjusted-specification estimates in Table~\ref{tab:stage1_composite} ($-0.076$ and $-0.066$ respectively), confirming that the standalone \texttt{Post} term makes little difference to the headline estimate; the covariate coefficients (young household head, female household head, ethnicity, disability, lone-parent status, family size) are all signed in the expected direction and highly significant, consistent with \citeauthor{andersen2025}'s own reported pattern.

\input{Tables/stage1_case_exact.tex}

\subsection{Item-level baseline difference-in-differences}
\label{sec:res-items}

Item-by-item estimates of Equation~\ref{eq:baseline-did} under the Adjusted specification, with Benjamini--Hochberg correction applied across the ten resulting tests, are reported in full in Appendix Table~\ref{tab:stage1_items}. Only two items survive correction: Warm coat ($-0.007$, BH $p=0.010$) and Celebrations ($-0.025$, BH $p=0.010$). Bed/bedroom has the largest point estimate ($-0.102$) but does not survive even at the raw 10\% level, and is estimated on a substantially thinner sample (around 7,000 observations, versus 45,000 or more for most other items), since it is only asked under the old, pre-FYE2024 question wording. The remaining seven items are not significant at any level.

\subsection{Item-stacked difference-in-differences}
\label{sec:res-stacked}

The item-stacked estimates from Section~\ref{sec:item-stacked} -- a pooled specification (one average effect across all ten items) and an item-interacted specification (one coefficient per item), both in Simple and Adjusted form -- are reported in full in Appendix Table~\ref{tab:stage2_stacked}. The pooled effect is small and not significant in either specification ($-0.015$, $p=0.153$ Simple; $-0.015$, $p=0.137$ Adjusted), and a joint Wald test of the null that all ten item-specific effects are zero is only marginal (Simple stat$=1.730$, $p=0.068$; Adjusted stat$=1.640$, $p=0.089$), clearing neither the conventional 5\% threshold.

Two items stand out in the item-interacted specification: Celebrations ($-0.015$ Simple, $-0.016$ Adjusted, both $p<0.05$) and, more marginally, Holiday away ($-0.060$ Simple, $-0.057$ Adjusted, both $p<0.10$). Comparing this table to Table~\ref{tab:stage1_items}'s independent-regression estimates reveals a genuine methodological divergence rather than a covariate mismatch, since both tables now share identical covariates: Celebrations is robust to the choice between the two approaches (independent: $-0.025$, BH-significant; stacked: $-0.016$, $p=0.043$), but Warm coat is not -- it survives Benjamini--Hochberg correction under independent regressions ($p=0.002$) yet is not significant in the pooled, household-clustered model ($p=0.264$). Since the only thing that differs between the two specifications is whether the ten items are fit as independent regressions or as one joint model accounting for within-child correlation across items, this divergence is a genuine open question about which is the more appropriate way to test a single item's effect when the same respondent answers all ten questions, rather than an artefact of specification choice, and is developed further in the discussion.

\subsection{Double/debiased machine learning difference-in-differences}
\label{sec:res-dml}

Table~\ref{tab:stage3_dml_composite} reports the ATET estimates for the official MDCH flag under every combination of nuisance model and covariate breadth (Section~\ref{sec:dml}), plus a lasso/random-forest ensemble for the lean specification only, since the ensemble did not complete within a feasible runtime for the wide specification.

\input{Tables/stage3_dml_composite.tex}

Together these allow the full specification ladder to be assembled: OLS-simple ($-0.062$) $\rightarrow$ OLS-adjusted ($-0.076$) $\rightarrow$ DML-lean-lasso ($-0.035$, not significant) $\rightarrow$ DML-lean-random-forest ($-0.040$, $p<0.10$) $\rightarrow$ DML-wide-lasso ($-0.028$, not significant) $\rightarrow$ DML-wide-random-forest ($-0.045$, $p<0.05$). Every specification points in the same direction and at a broadly similar magnitude ($-0.028$ to $-0.076$); lasso is consistently the weakest cell at both covariate breadths, consistent with linear nuisance models understating the effect relative to random forest, while the two random-forest cells are close enough to each other that covariate breadth looks secondary to estimator flexibility here. OLS's adjustment step overstates the effect relative to every DML cell, suggesting the covariate-adjusted OLS estimate in Section~\ref{sec:res-baseline} partly reflects functional-form misspecification rather than the transfer's full causal effect. A lasso/random-forest ensemble estimated for the lean specification ($-0.041$, $p<0.10$) sits between the two individual lean-covariate methods, as expected.

\subsection{Item-level and stacked machine-learning DiD}
\label{sec:res-item-dml}

Item-level DML estimates (Section~\ref{sec:item-dml}), using the wide covariate set with lasso and random forest nuisance models, are reported in full in Appendix Table~\ref{tab:stage4_dml_items}. Point estimates broadly track the item-stacked OLS-adjusted estimates in Appendix Table~\ref{tab:stage2_stacked}: Celebrations' lasso estimate ($-0.016$) matches OLS almost exactly, and Warm coat's ($-0.007$ lasso, $-0.008$ random forest) sits close to both OLS specifications. But standard errors are wide enough that even these two items, which survive Benjamini--Hochberg correction under independent baseline OLS (Section~\ref{sec:res-items}), fail to clear correction here: none of the ten items is BH-significant under either lasso or random forest. Holiday away comes closest to conventional significance under lasso alone ($-0.056$, raw $p=0.078$), but does not survive correction either.

This is extended in Appendix Table~\ref{tab:stage5_stacked_ml} to a single, jointly estimated model across items, sharing cross-fitted propensity models across items and predicting all item outcomes at once via a multivariate outcome-regression learner (Section~\ref{sec:item-dml}). Bed/bedroom and Playgroup attendance are dropped from this stage, since requiring one consistent complete-case sample across every included item leaves too thin a joint sample if either is retained, leaving eight items on a smaller joint sample than any single item's own Appendix Table~\ref{tab:stage4_dml_items} sample. None of the eight items is BH-significant (all $p_{BH}\geq0.371$): Warm coat, School trip and Outdoor play area tie closest.

The more informative pattern is the direction of change relative to the independently fitted lasso estimates in Table~\ref{tab:stage4_dml_items}: Holiday away, the strongest individual-item candidate under independent fitting, weakens sharply under stacking ($-0.056\to-0.023$), while a few other items (Warm coat, Outdoor play area, School trip) move modestly the other way, still nowhere near significance. This instability -- which item looks closest to significance depends on whether items are modelled independently or jointly -- is itself evidence against a genuine item-level effect rather than for one. Taken together with Appendix Table~\ref{tab:stage1_items}'s independent-OLS estimates, the item-level null finding now holds under three independent estimation strategies (independent lasso, independent random forest, and joint stacked lasso): only two items ever survive multiple-testing correction, and both do so only under simple OLS, not under any form of doubly robust re-estimation.

\subsection{Falsification checks: Wales and Northern Ireland}
\label{sec:res-placebo}

As a falsification check on the Scotland--England design, the headline specification is re-estimated with Wales, and separately Northern Ireland, as the treated unit against England as control, using SCP's actual rollout date as a pseudo-treatment date; since neither nation received SCP, a null coefficient is the reassuring result. Table~\ref{tab:placebo_wales_ni} reports the results for all four outcomes.

\input{Tables/placebo_wales_ni.tex}

Wales passes cleanly on all three (coefficients close to zero, $p$-values between 0.26 and 0.98 in both the Simple and Adjusted specifications). Northern Ireland passes on two of the three (official MDCH flag and severe deprivation, both null), but shows a significant, positive coefficient on any deprivation specifically (Simple $0.086$, $p=0.035$; Adjusted $0.080$, $p=0.036$), indicating that measured deprivation rose in Northern Ireland relative to England around SCP's rollout date on this one composite. This coefficient is of the opposite sign to Scotland's own (negative) effect, so it does not mimic or contaminate the headline estimate; any deprivation is also the composite that shows the least conservative behaviour elsewhere in this paper (Section~\ref{sec:parallel-trends}), consistent with roughly one false positive being unsurprising across the 12 comparator$\times$outcome$\times$specification cells tested here. One candidate real-world explanation -- a change to Northern Ireland's own welfare "mitigations" scheme -- was checked and ruled out, since that scheme was extended, not curtailed, over this period. The cause of the Northern Ireland/any-deprivation anomaly is left as an open, flagged limitation rather than resolved.

\subsection{Robustness checks}
\label{sec:res-robustness}

Three further checks support the baseline design, alongside the Wales/Northern Ireland falsification check reported above (Section~\ref{sec:res-placebo}). First, excluding FYE2022 (the partial-rollout year, in which under-6 Scottish children were already receiving SCP for part of the year) from the pre-period leaves the headline result unchanged, confirming it is not driven by the transition year itself. Second, the official MDCH flag is confirmed to be computed on the same basis in Scotland and England, with no region-specific necessities weighted differently between the two nations, ruling out a measurement-consistency confound. Third, since Scotland also runs other Scotland-only, child-targeted transfers alongside SCP (Section~\ref{sec:identification-assumptions}), the Adjusted specification is re-estimated on children aged 6--15 only, excluding the under-6 band where Best Start Grant and Best Start Foods eligibility could be entangled with SCP's own effect. Despite losing roughly a third of the sample, every coefficient stays consistent in sign and similar in magnitude to the full-sample estimates -- if anything larger for the official MDCH flag ($-0.087$, $p<0.05$, versus $-0.076$, $p<0.01$) -- reassuring evidence against Best Start driving the headline result through the under-6 cohort specifically; full results across all five outcomes are reported in Appendix Table~\ref{tab:stage1_age_restricted}.
